% ======================================================================
%  Documento técnico: flujo_sync_faq_flujo.py
%  Explicación detallada del mecanismo de sincronización automática
%  FAQ markdown -> diagramas de flujo (.tex/.pdf)
%  Estilo infográfico replicado de docs/AGENTE_IA/*.tex
% ======================================================================
\documentclass[11pt,a4paper]{article}

% -- Codificación, fuentes y lenguaje -----------------------------------
\usepackage[utf8]{inputenc}
\usepackage[T1]{fontenc}
\usepackage[default,scale=0.95]{sourcesanspro}
\renewcommand{\familydefault}{\sfdefault}
\usepackage[spanish,es-noshorthands,es-tabla]{babel}

% -- Geometría y espaciado ----------------------------------------------
\usepackage[top=2.2cm, bottom=2.5cm, left=2.2cm, right=2.2cm, headheight=15pt]{geometry}
\usepackage{setspace}
\setstretch{1.15}
\usepackage{parskip}

% -- Matemáticas --------------------------------------------------------
\usepackage{amsmath}
\usepackage{amssymb}

% -- TikZ (símbolos ISO 5807 / ANSI X3.5) -------------------------------
\usepackage{tikz}
\usetikzlibrary{babel}
\usetikzlibrary{arrows.meta, positioning, shapes.geometric, shapes.symbols,
                decorations.pathmorphing, calc}

% -- Tablas y listas ----------------------------------------------------
\usepackage{booktabs}
\usepackage{tabularx}
\usepackage{array}
\usepackage{enumitem}

% -- Código fuente ------------------------------------------------------
\usepackage{listings}
\lstset{
  basicstyle=\ttfamily\footnotesize,
  breaklines=true,
  showstringspaces=false,
  frame=single,
  rulecolor=\color{grisLinea},
  backgroundcolor=\color{grisPapel},
  language=Python,
  literate=
    {á}{{\'a}}1 {é}{{\'e}}1 {í}{{\'i}}1 {ó}{{\'o}}1 {ú}{{\'u}}1
    {Á}{{\'A}}1 {É}{{\'E}}1 {Í}{{\'I}}1 {Ó}{{\'O}}1 {Ú}{{\'U}}1
    {ñ}{{\~n}}1 {Ñ}{{\~N}}1 {ü}{{\"u}}1 {¿}{{?`}}1 {¡}{{!`}}1,
}

% -- Rutas de archivo con quiebre de línea ------------------------------
\usepackage{url}

% -- Colores, cajas e iconos ---------------------------------------------
\usepackage[dvipsnames]{xcolor}
\usepackage{tcolorbox}
\tcbuselibrary{skins, breakable}
\usepackage{fontawesome5}

% -- Secciones con barra de color (infografía) ---------------------------
\usepackage{titlesec}
\titleformat{\section}
  {\normalfont\LARGE\bfseries\color{azulNoche}}
  {\colorbox{cyanNeon}{\color{fondoOscuro}\thesection}}{0.7em}{}
  [\vspace{2pt}\color{cyanNeon}\rule{\linewidth}{1.8pt}]
\titlespacing*{\section}{0pt}{24pt}{10pt}
\titleformat{\subsection}
  {\normalfont\large\bfseries\color{azulNoche}}
  {\textcolor{cyanNeon}{\thesubsection}}{0.7em}{}
  [\vspace{1pt}\color{grisLinea}\rule{\linewidth}{0.6pt}]
\titlespacing*{\subsection}{0pt}{14pt}{6pt}

% -- Cabeceras y pies de página ------------------------------------------
\usepackage{fancyhdr}
\usepackage{lastpage}
\pagestyle{fancy}
\fancyhf{}
\renewcommand{\headrulewidth}{0pt}
\renewcommand{\footrulewidth}{0pt}
\fancyfoot[C]{%
  \begin{minipage}{\linewidth}
    \vspace{2pt}
    \color{cyanNeon}\rule{\linewidth}{1.2pt}\\[2pt]
    \centering\color{grisTexto}\footnotesize
    Página \thepage\ de \pageref{LastPage}
  \end{minipage}%
}

% -- Hiperenlaces ---------------------------------------------------------
\usepackage[colorlinks=true, linkcolor=azulNoche, urlcolor=cyanNeon]{hyperref}
\sloppy
\emergencystretch 3em

% -- Paleta de colores ----------------------------------------------------
\definecolor{fondoOscuro}{HTML}{0D1B2A}
\definecolor{verdeTurquesa}{HTML}{004D40}
\definecolor{azulNoche}{HTML}{1B3A6B}
\definecolor{azulMedio}{HTML}{2A5298}
\definecolor{cyanNeon}{HTML}{00C8E0}
\definecolor{verdeSignal}{HTML}{00B887}
\definecolor{naranjaVivo}{HTML}{FF7043}
\definecolor{rojoAlerta}{HTML}{E53935}
\definecolor{grisPapel}{HTML}{F4F6FA}
\definecolor{grisLinea}{HTML}{C8D0E0}
\definecolor{grisTexto}{HTML}{6B7A99}
\definecolor{amarilloNota}{HTML}{FFB300}

% -- Banda de título ------------------------------------------------------
\newcommand{\iconoBanda}{sync-alt}
\newcommand{\bandaTitulo}[3][fondoOscuro]{%
  \begin{tcolorbox}[
    enhanced,
    colback=#1, colframe=#1,
    boxrule=0pt, arc=8pt, outer arc=8pt,
    width=\linewidth,
    top=18pt, bottom=6pt, left=14pt, right=14pt,
    borderline south={3.5pt}{0pt}{cyanNeon},
  ]
    \begin{minipage}[c]{0.07\linewidth}
      \centering
      \textcolor{cyanNeon}{\fontsize{30}{30}\selectfont\faIcon{\iconoBanda}}
    \end{minipage}%
    \hspace{8pt}%
    \begin{minipage}[c]{0.88\linewidth}
      {\fontsize{20}{24}\selectfont\bfseries\color{white}#2}\par\vspace{5pt}%
      \textcolor{cyanNeon!80!white}{\small\faIcon{angle-right}~#3}%
    \end{minipage}
    \vspace{4pt}
  \end{tcolorbox}%
  \vspace{8pt}%
}

% -- Tarjetas --------------------------------------------------------------
\tcbset{
  cajaTitulo/.style={
    enhanced, breakable,
    arc=6pt, outer arc=6pt,
    colback=azulNoche!10!grisPapel, colframe=azulNoche,
    boxrule=1pt,
    fonttitle=\bfseries, coltitle=white,
    attach boxed title to top left={yshift=-3mm, xshift=6mm},
    boxed title style={colback=azulNoche, colframe=azulNoche, arc=4pt, boxrule=0pt},
    drop fuzzy shadow=azulNoche!25!white,
    top=8pt, bottom=8pt, left=10pt, right=10pt,
  },
  cajaContenido/.style={
    enhanced, breakable,
    arc=5pt, outer arc=5pt,
    colback=grisPapel, colframe=grisLinea,
    boxrule=0.8pt,
    fonttitle=\bfseries\small, coltitle=azulNoche,
    top=6pt, bottom=6pt, left=10pt, right=10pt,
  },
  cajaConcepto/.style={
    enhanced, breakable,
    arc=6pt, outer arc=6pt,
    colback=verdeSignal!8!white,
    colframe=verdeSignal, boxrule=1pt,
    borderline west={4pt}{0pt}{verdeSignal},
    fonttitle=\bfseries\small\color{white},
    title={\faIcon{lightbulb}~Concepto clave},
    attach boxed title to top left={yshift=-3mm, xshift=6mm},
    boxed title style={colback=verdeSignal!80!black, arc=4pt, boxrule=0pt},
    drop fuzzy shadow=verdeSignal!20!white,
    left=10pt, right=10pt, top=8pt, bottom=8pt,
  },
  cajaMejora/.style={
    enhanced, breakable,
    arc=6pt, outer arc=6pt,
    colback=naranjaVivo!8!white, colframe=naranjaVivo,
    boxrule=1pt,
    borderline west={4pt}{0pt}{naranjaVivo},
    fonttitle=\bfseries\small, coltitle=naranjaVivo!80!black,
    drop fuzzy shadow=naranjaVivo!20!white,
    top=6pt, bottom=6pt, left=10pt, right=10pt,
  },
}

\newcommand{\iconotexto}[2]{\textcolor{cyanNeon}{\faIcon{#1}}~#2}

% -- Estilos de flujo (ISO 5807 / ANSI X3.5) ------------------------------
\tikzset{
  flujo/.style={
    >=Stealth,
    font=\footnotesize\sffamily,
  },
  terminador/.style={
    draw=azulNoche, fill=azulNoche!8!white, rounded corners=3pt,
    text width=1.6cm, align=center, inner sep=2mm,
    font=\footnotesize\sffamily,
  },
  proceso/.style={
    draw=azulNoche, fill=cyanNeon!10!white,
    text width=1.6cm, align=center, inner sep=2mm,
    font=\footnotesize\sffamily,
  },
  entradasalida/.style={
    draw=azulNoche, fill=verdeSignal!10!white,
    trapezium, trapezium left angle=75, trapezium right angle=105,
    text width=1.8cm, align=center, inner sep=2mm,
    font=\footnotesize\sffamily,
  },
  decision/.style={
    draw=azulNoche, fill=amarilloNota!20!white,
    diamond, aspect=1.6, text width=1.5cm, align=center, inner sep=1pt,
    font=\footnotesize\sffamily,
  },
  almacenamiento/.style={
    draw=azulNoche, fill=grisPapel,
    cylinder, shape border rotate=90, aspect=0.3,
    text width=1.7cm, align=center, inner sep=2mm,
    font=\footnotesize\sffamily,
  },
  documento/.style={
    draw=azulNoche, fill=azulMedio!8!white,
    text width=1.8cm, align=center, inner sep=2mm,
    font=\footnotesize\sffamily,
  },
  nota/.style={
    draw=grisLinea, dashed, fill=grisPapel,
    text width=1.8cm, align=center, inner sep=2mm,
    font=\scriptsize\sffamily,
  },
  etiqueta/.style={font=\scriptsize\sffamily, fill=white, fill opacity=0.75,
                   text opacity=1, inner sep=1pt},
}

% -- Cabecera específica del documento -------------------------------------
\fancyhead[L]{\small\color{azulNoche}\textbf{ELECTRÓNICA} --- Sincronización automática FAQ$\rightarrow$Flujo}
\fancyhead[R]{\small\color{grisTexto}12/09/2026}
\renewcommand{\headrulewidth}{0pt}

% ======================================================================
\begin{document}

\bandaTitulo{Sincronización Automática FAQ $\rightarrow$ Diagramas de Flujo}{Ficha técnica del mecanismo de 3 capas que mantiene \texttt{faq\_higiene\_estado\_sesion\_flujo.\{tex,pdf\}} alineados con \texttt{faq\_higiene\_estado\_sesion.md}, sin intervención manual}

% ======================================================================
\section{Objetivo y motivación}

El documento \texttt{docs/AGENTE\_IA/faq\_higiene\_estado\_sesion\_flujo.tex} es una
\textbf{traducción semántica} (no mecánica ni 1:1) del FAQ
\texttt{faq\_higiene\_estado\_sesion.md}. Los cinco diagramas ISO 5807 que contiene
resumen la cadena de integridad, la guardia MCP por \texttt{mtime}, los veredictos de
\texttt{estado\_sesion.py check/clean} y el eslabón semántico de \texttt{bitacoras.py}.

El problema que motiva este flujo es de \emph{consistencia}: si alguien edita el FAQ
(añade un veredicto, cambia una fecha, corrige una frase), el diagrama queda obsoleto.
Hacer la traducción a mano es lento, y la traducción ``desde cero'' por un LLM rompe la
geometría de los diagramas (coordenadas, estilos y número de nodos son frágiles).

\begin{tcolorbox}[cajaConcepto]
\faIcon{bullseye}~Objetivo: cuando se edite el \texttt{.md}, el flujo detecta el cambio
por \textbf{hash SHA-256}, decide \emph{qué} hay que traducir, lo traduce de forma
\textbf{híbrida} (campos deterministas por script, semántica vía LLM quirúrgico),
recompila el PDF y \textbf{verifica} el resultado sin necesidad de visión humana.
\end{tcolorbox}

El resultado es un bucle de auto-curación determinista: la misma entrada (mismo \texttt{.md}
en el mismo estado) produce siempre el mismo comportamiento, y la única parte probabilística
(el LLM) se mantiene acotada y \emph{validada} antes de tocar el documento.

% ======================================================================
\section{Arquitectura de 3 capas}

El proyecto ELECTRONICA impone una separación estricta de responsabilidades. Este flujo la
respeta al pie de la letra:

\begin{table}[h!]
\centering
\small
\begin{tabularx}{\linewidth}{@{}>{\bfseries}p{2.4cm}>{\raggedright\arraybackslash}p{3.2cm}>{\raggedright\arraybackslash}X@{}}
\toprule
Capa & Artefacto & Rol (qu\'e / c\'omo) \\
\midrule
1. Directiva & \path{directives/sync_faq_a_flujo.yaml} & SOP de nivel medio: qu\'e hacer, paso a paso, y qu\'e hacer ante cada caso l\'imite (edge case). \\
2. Orquestaci\'on & \path{flujo_sync_faq_flujo.py} & Toma de decisiones: comparar hash, planear, encadenar ejecuci\'on, avisar de coste, compilar, copiar, alertar. \\
3. Ejecuci\'on & \path{execution/regenerar_faq_flujo.py} \newline \path{execution/verificar_pdf.py} & C\'omo: scripts deterministas con CLI, salidas JSON y c\'odigos de salida categorizados. \\
\bottomrule
\end{tabularx}
\end{table}

A esta tr\'iada se suma un cambio \emph{aditivo} en la pieza compartida
\texttt{execution/compile\_latex.py}: un par\'ametro nuevo \texttt{clean=False} que permite al
orquestador conservar el log compilaci\'on hasta despu\'es de la verificaci\'on (ver
Secci\'on~\ref{sec:compile}).

\begin{tcolorbox}[cajaMejora]
\faIcon{handshake}~\textbf{Reutilización antes que construcción.} El flujo no reimplementa
nada: reutiliza \texttt{execution/enrutador.py} (decisión de tier), \texttt{execution/llm\_client.py}
(chat por OpenRouter), \texttt{execution/compile\_latex.py} (compilación), \texttt{execution/alert\_user.py}
(notificación audible) y los patrones JSON ya documentados en \texttt{.agent/python.md}.
\end{tcolorbox}

% ======================================================================
\section{Visión general del flujo completo}

\subsection{La cadena, de principio a fin}

Cada ejecución del orquestador recorre esta secuencia (fases E1--E5 del SOP):

\begin{center}
\begin{tikzpicture}[flujo, scale=0.66, transform shape, x=1cm, y=1cm]
  % Nodo inicial
  \node[terminador] (inicio) at (0,0) {Orquesta\\ \texttt{--once}};
  \node[proceso] (f1) at (0,-1.6) {SHA-256\\ del .md};
  \node[almacenamiento, text width=3cm, align=center] (estado) at (3,-1.6) {Estado\\ \texttt{faq\_flujo\_sync.json}};
  \node[decision] (f2) at (0,-3.3) {¿Cambió?};
  \node[terminador] (fin0) at (3,-3.3) {Fin\\ sin cambios};
  \node[proceso] (p1) at (0,-5.0) {Clasificar plan};
  \node[decision] (f3) at (0,-6.7) {¿Cambio\\ estructural?};
  \node[proceso] (det) at (-3,-6.7) {Inyectar\\ campos};
  \node[proceso] (llm1) at (0,-8.4) {LLM quirúrgico\\ (enrutador)};
  \node[nota] (notallm) at (3.2,-7.6) {Aviso previo\\ de créditos\\ OpenRouter};
  \node[proceso] (c1) at (0,-10.2) {Reescribir .tex};
  \node[proceso] (c2) at (0,-11.9) {pdflatex $\times$\\ 2 pasadas};
  \node[proceso] (c3) at (0,-13.6) {Copiar PDF};
  \node[proceso] (c4) at (0,-15.3) {Verificar\\ PDF};
  \node[proceso] (c5) at (0,-17.0) {Estado +\\ snapshot};
  \node[terminador] (fin) at (0,-18.7) {Fin\\ sincronizado};

  % Aristas
  \draw[->] (inicio.south) -- (f1.north);
  \draw[<->, dashed] (f1.east) -- (estado.west);
  \draw[->] (f1.south) -- (f2.north);
  \draw[->] (f2.east) -- node[etiqueta, right, pos=0.5]{No} (fin0.west);
  \draw[->] (f2.south) -- node[etiqueta, left, pos=0.5]{Sí} (p1.north);
  \draw[->] (p1.south) -- (f3.north);
  \draw[->] (f3.west) -- node[etiqueta, left, pos=0.5]{No} (det.east);
  \draw[->] (f3.south) -- node[etiqueta, left, pos=0.5]{Sí} (llm1.north);
  \draw[dashed] (llm1.east) -- (notallm.west);
  \draw[->] (det.south) |- (c1.west);
  \draw[->] (llm1.south) -- (c1.north);
  \draw[->] (c1.south) -- (c2.north);
  \draw[->] (c2.south) -- (c3.north);
  \draw[->] (c3.south) -- (c4.north);
  \draw[->] (c4.south) -- (c5.north);
  \draw[->] (c5.south) -- (fin.north);
\end{tikzpicture}
\end{center}

\emph{Nota: la decisión central ``¿cambio estructural?'' la toma \texttt{regenerar\_faq\_flujo.py}
comparando el \texttt{.md} con su snapshot (Sección~\ref{sec:clasific}), no el chat del agente.}

\subsection{Geografía de archivos}

\begin{table}[h!]
\centering
\small
\begin{tabularx}{\linewidth}{@{}>{\raggedright\arraybackslash}p{8.2cm}>{\raggedright\arraybackslash}X@{}}
\toprule
Archivo & Papel \\
\midrule
\path{docs/AGENTE_IA/faq_higiene_estado_sesion.md} & Fuente de verdad del FAQ (se edita). \\
\path{docs/AGENTE_IA/faq_higiene_estado_sesion_flujo.tex} & Plantilla de diagramas; se edita \emph{in situ}; nunca se genera desde cero. \\
\path{docs/AGENTE_IA/faq_higiene_estado_sesion_flujo.pdf} & Entregable compilado (entrada del usuario). \\
\path{.tmp/faq_flujo_sync.json} & Estado del flujo: hashes md/tex, marca temporal, v\'ia, secciones, edits. \\
\path{.tmp/faq_flujo_md_snapshot.md} & Copia del \texttt{.md} en el \'ultimo arranque (base del diff de clasificaci\'on). \\
\path{.tmp/faq_flujo_sync.log} & Log append-only (hora UTC) de cada corrida. \\
\path{.tmp/latex_build/} & Directorio de compilaci\'on (auxiliares + \texttt{.pdf} de build). \\
\bottomrule
\end{tabularx}
\end{table}

El estado en \texttt{.tmp/} es \emph{derivado}: cualquiera de esos ficheros puede eliminarse y
el flujo lo regenera en el siguiente arranque (la l\'inea base se establece sola).

% ======================================================================
\section{Detección de cambios por hash}

\label{sec:hash}
El disparador no depende de \texttt{inotify} ni de watchers del RAG: es un \textbf{polling por
hash} que también funciona con una sola ejecución puntual.

\begin{lstlisting}
def _sha256(path):
    return hashlib.sha256(path.read_bytes()).hexdigest()

estado = _leer_estado()                      # .tmp/faq_flujo_sync.json
md_sha = _sha256(md)
if not args.force and estado.get("md_sha256") == md_sha:
    return {..., "cambio": False, "message": "Sin cambios en el markdown (hash idén-tico)."}
\end{lstlisting}

\begin{itemize}[leftmargin=1.4em]
  \item \texttt{--force} salta la comparación y fuerza la cadena completa (equivalentemente,
        la primera ejecución sin estado previo también entra).
  \item \texttt{--watch [--interval N]} repite la comprobación cada $N$ segundos (mínimo 5) en
        un bucle hasta Ctrl-C; cada corrida imprime su propio JSON y log.
  \item El hash también se guarda en el estado tras sincronizar, de modo que la siguiente
        corrida detecta \emph{el próximo} cambio (no el mismo).
\end{itemize}

\begin{tcolorbox}[cajaConcepto]
\faIcon{balance-scale}~La detección por hash es \textbf{determinista} (función pura de los
bytes del archivo) y \textbf{barata}: sin cambio no hay llm, ni compilación, ni créditos.
\end{tcolorbox}

% ======================================================================
\section{Clasificación del cambio: determinista vs. semántico}

\label{sec:clasific}
Antes de llamar a ningún LLM, el script \texttt{regenerar\_faq\_flujo.py} responde a la
pregunta: \emph{¿este cambio necesita traducción semántica o basta con una sustitución?}

\subsection{Campos conocidos (inyección determinista)}

Hay tres cosas del FAQ que se traducen por \textbf{regex exactas}, sin LLM:

\begin{table}[h!]
\centering
\small
\begin{tabularx}{\linewidth}{@{}p{2.9cm}p{5.4cm}X@{}}
\toprule
Campo & Regex & Efecto en el .tex \\
\midrule
Fecha de la plantilla & \texttt{FECHA\_PLANTILLA\ (YYYY-MM-DD)} & Actualiza la fecha ISO y el encabezado día/mes/año (DD/MM/AAAA). \\
Modelo obsoleto & \texttt{deepseek-[\textbackslash{}w.-]+} & Sustituye el identificador del modelo (\texttt{deepseek-v4-pro} $\rightarrow$ el vigente). \\
Bitácoras antiguas & \texttt{\textbackslash{}d+\textbackslash{}s+antiguas} & Ajusta el número ``N antiguas'' en la sección de bitácoras. \\
\bottomrule
\end{tabularx}
\end{table}

Estas dejan de tocar el LLM y devuelven \texttt{via: "determinista"}. El script aplica la
sustitución en memoria y solo la escribe si hubo cambio real (\texttt{tex\_cambiado}).

\subsection{El diff contra el snapshot}

Para lo demás se usa \texttt{difflib.SequenceMatcher} entre el \texttt{.md} actual y la copia
\texttt{.tmp/faq\_flujo\_md\_snapshot.md} del último arranque. Los \emph{opcodes} de cambio se
agrupan por sección del FAQ:

\begin{itemize}[leftmargin=1.4em]
  \item Las cabeceras \texttt{\#\# N. ...} se indexan con una regex y cada línea del diff se
        asigna a su sección (\texttt{\_seccion\_de\_linea}).
  \item Una sección se marca \texttt{semantico=True} si alguna línea cambia después de
        \emph{normalizar} los campos conocidos (\texttt{\_scrub\_line} reemplaza fecha/modelo/N
        por marcadores de posición), o si el número de líneas viejo $\neq$ nuevo.
  \item Las secciones ``front'' (metadatos) y ``Referencias cruzadas'' se marcan ignoradas.
\end{itemize}

\begin{lstlisting}
def _scrub_line(line):
    line = re.sub(_RE_FECHA, "<FECHA>", line)
    line = re.sub(_RE_MODELO, "<MODELO>", line)
    line = re.sub(_RE_ANTIGUAS, "<N> antiguas", line)
    return line

for o, n in zip(old_lines, new_lines):
    if _scrub_line(o) != _scrub_line(n):
        entry["semantico"] = True        # cambio real de contenido
\end{lstlisting}

\subsection{Mapeo sección del FAQ $\rightarrow$ bloque del .tex}

El \texttt{.tex} se divide en 4 bloques usando marcadores de comentario propios
(\texttt{\%} seguido de dos líneas de caja, \texttt{U+2550}, y ``Sección N:''), El bloque 4
termina en \texttt{\textbackslash end\{document\}}:

\begin{lstlisting}
_MARKERS = {n: "% \u2550\u2550 Secci\u00f3n " + str(n) + ":" for n in (1,2,3,4)}
def _blocks_tex(tex):
    pos = {n: tex.find(S) for n, S in _MARKERS.items()}
    end = tex.find("\\end{document}")
    return {1: (pos[1], pos[2]), 2: (pos[2], pos[3]),
            3: (pos[3], pos[4]), 4: (pos[4], end)}
\end{lstlisting}

Así, un cambio en la sección 2 del FAQ solo dispara evaluación sobre el bloque 2 del \texttt{.tex};
el resto del documento queda intacto por construcción. Esta es una divergencia \emph{consciente}
entre la numeración de Markdown y la de bloques: los diagramas no siguen 1:1 los encabezados.

\subsection{La salida del plan}

Con \texttt{--plan} el script no llama al LLM ni escribe: devuelve el mapa de decisión. Ejemplo
real tras un cambio en la sección 2:

\begin{lstlisting}
{ "status": "ok", "tex_cambiado": false, "via": "plan",
  "known_fields": {},
  "cambios": [ {"seccion": "2", "via": "llm", "hunks": 1} ],
  "llm_necesario": true,
  "nota": "La sincronizacion real invocara el LLM via OpenRouter (consume creditos)." }
\end{lstlisting}

El orquestador usa este veredicto para \textbf{avisar al usuario del coste antes} de consumir
créditos, y para abortar limpiamente si se pasa \texttt{--no-llm} y el cambio es estructural.

% ======================================================================
\section{Re-traducción semántica (LLM quirúrgico)}

\label{sec:llmquirc}

Cuando hay cambios estructurales, el flujo no regenera el bloque: lo \emph{edita por
parches}. El LLM recibe como plantilla el bloque \texttt{.tex} actual y produce un JSON de
edits; el script valida y aplica. Así se preserva la geometría dibujada y ajustada
manualmente.

\subsection{Routing determinista (sin decisión en el chat)}

\label{sec:routing}
El tier se decide con \texttt{execution/enrutador.py} a partir del descriptor
\texttt{--task conversion} y los \textbf{tokens medidos} del corpus (markdown afectado +
bloque \texttt{.tex} + inventario del diff):

\begin{lstlisting}
task = "conversion"
decision = decide(task, tokens, critico, False, None)   # funcion pura
candidates = [decision["model"]] + (decision.get("fallback") or [])[:2]
if model:                                                # --modelo explicito
    candidates = [model] + [c for c in candidates if c != model]
\end{lstlisting}

\begin{itemize}[leftmargin=1.4em]
  \item Por defecto el tier es \texttt{flash} (\texttt{google/gemini-3.7-flash}, barato y
        suficiente para edición de tablas/texto).
  \item \texttt{--critico} escala a \texttt{opus}; el fallback cost-aware de la cadena
        (\texttt{flash $\rightarrow$ deepseek $\rightarrow$ glm}) solo se usa ante 429/errores,
        con retry budget máximo 3 intentos.
  \item Cada decisión se registra en \texttt{.tmp/routing\_log.jsonl} (telemetría), y la
        llamada consume créditos de \texttt{OPENROUTER\_API\_KEY} vía
        \texttt{execution/llm\_client.openrouter\_chat}.
\end{itemize}

\subsection{El prompt de editor quirúrgico}

El \texttt{system} prompt fija reglas invariables; la más importante es \emph{no tocar la
geometría}. El \texttt{user} prompt entrega el cambio detectado (fragmentos \texttt{ANTES/
AHORA}) y el bloque actual entre marcadores.

\begin{lstlisting}
system = (
 "Eres un editor quirurgico de codigo LaTeX ...\n"
 "1. NO modifiques la geometria: ni coordenadas, ni estilos \node[...], ni el numero de "
 "nodos, ni las etiquetas internas (T1, P1, D1, E1, A1, N1...).\n"
 "2. Solo edita textos: filas de la tabla 'Leyenda de etiquetas', parrafos explicativos...\n"
 "3. Puedes anadir filas a la tabla (minimas) o reescribir significados. "
 "PROHIBIDO anadir nodos o secciones.\n"
 "4. Manten el LaTeX valido: en tablas escapa & como \& y _ como \_ dentro de \texttt.\n"
 "5. Devuelve SOLO JSON: {\"edits\":[{\"old\":\"subcadena unica\",\"new\":\"texto\"}]}. "
 "Si no hay nada que cambiar, {\"edits\":[]}.\n"
 "6. No escribas nada fuera del JSON."
)
\end{lstlisting}

\subsection{De la respuesta cruda a edits seguros}

\begin{enumerate}[leftmargin=1.6em]
  \item \texttt{\_parse\_edits}: extrae el primer objeto JSON balanceado
        (\texttt{\_find\_balanced\_json}, que respeta strings y escapes -- ver
        \texttt{.agent/python.md}), elimina trailing commas y hace \texttt{json.loads}.
  \item \texttt{\_edits\_validos}: exige que cada \texttt{old} sea una subcadena \textbf{única}
        del bloque y que el \texttt{new} no contenga tokens estructurales prohibidos
        (\texttt{\textbackslash begin\{tikzpicture\}}, \texttt{\textbackslash section},
        \texttt{\textbackslash documentclass}, \dots). Cualquier incumplimiento aborta con error.
  \item \texttt{\_aplicar\_edits}: sustitución estricta (re-valida unicidad antes de cada
        \texttt{replace}).
  \item \textbf{Invariante geométrico}: tras aplicar todos los edits de un bloque, se cuenta
        \texttt{\textbackslash node[} antes y después; si difieren, el script aborta \emph{sin
        escribir} (código 5).
\end{enumerate}

\begin{lstlisting}
n_antes   = block.count("\\node[")
n_despues = nuevo_block.count("\\node[")
if n_antes != n_despues:
    raise RuntimeError("El LLM altero la geometria del bloque ... Abortando sin escribir.")
\end{lstlisting}

\subsection{Escritura atómica}

Solo después de validar todos los bloques se escribe el fichero, y se hace de forma atómica
(archivo temporal + \texttt{os.replace}) para no dejar nunca un \texttt{.tex} a medias:

\begin{lstlisting}
if not (args.dry_run or not resultado["tex_cambiado"]):
    fd, tmp = tempfile.mkstemp(dir=str(tex_path.parent), suffix=".tmp")
    with os.fdopen(fd, "w", encoding="utf-8") as fh:
        fh.write(tex)
    os.replace(tmp, tex_path)
\end{lstlisting}

Un LLM que devuelva \texttt{\{"edits":[]\}} (porque el bloque ya está al día) es un
\emph{resultado correcto}: la cadena valida el JSON vacío y continúa con
\texttt{tex\_cambiado: False}, sin recompilar.

% ======================================================================
\section{Compilación y el parámetro \texttt{clean=False}}

\label{sec:compile}
El orquestador compila el \texttt{.tex} regenerado con la pieza compartida
\texttt{compile\_latex\_code} (2 pasadas de \texttt{pdflatex} para resolver referencias) al
directorio \texttt{.tmp/latex\_build/}:

\begin{lstlisting}
res = compile_latex_code(contenido, job_name=JOB_NAME,
                         output_dir=str(BUILD_DIR), clean=False)
...
pdf_build = BUILD_DIR / (JOB_NAME + ".pdf")
shutil.copyfile(pdf_build, pdf)        # publica el entregable en docs/AGENTE_IA/
\end{lstlisting}

El parámetro nuevo \texttt{clean=False} es el único cambio tocado en la pieza compartida. ¿Por
qué? Por defecto \texttt{compile\_latex\_code} borra los auxiliares (incluido el \texttt{.log})
al terminar -- perfecto para el MCP de LaTeX, que solo necesita el PDF. Pero el guardrail de
verificación necesita leer el \texttt{.log} (errores \texttt{!} y Overfull). Con
\texttt{clean=False} el orquestador conserva el log, ejecuta \texttt{verificar\_pdf.py},
y \textbf{luego} limpia él mismo con \texttt{clean\_latex\_aux\_files(BUILD\_DIR)}.

\begin{tcolorbox}[cajaMejora]
\faIcon{plug}~\textbf{Compatibilidad asegurada:} el cambio es aditivo con valor por defecto
\texttt{True}. Ningún llamante existente (\texttt{mcp\_latex\_server.py} y los demás flujos)
cambia de comportamiento.
\end{tcolorbox}

% ======================================================================
\section{Guardrail de verificación (sin visión humana)}

\label{sec:verificar}
\texttt{execution/verificar\_pdf.py} es la validación determinista post-compilación. Tiene
tres frentes:

\begin{enumerate}[leftmargin=1.6em]
  \item \textbf{Errores de compilación}: líneas del \texttt{.log} que empiezan por \texttt{!}.
        Su presencia vuelve la salida \texttt{status: "error"} y código 1.
  \item \textbf{Overfull \textbackslash hbox}: se reportan (informativos, cosméticos) pero no
        bloquean.
  \item \textbf{Solapes de texto}: parsea \texttt{pdftotext -bbox} (caja por palabra) y
        detecta pares de palabras cuyos rectángulos se solapan más de un umbral.
\end{enumerate}

\subsection{Cálculo del solape}

\begin{lstlisting}
SOLAPE_FRACCION = 0.35

def _solapa(r1, r2):
    dx = min(x2,u2) - max(x1,u1); dy = min(y2,v2) - max(y1,v1)
    if dx <= 0 or dy <= 0: return 0.0
    inter = dx * dy
    menor = min(area(r1), area(r2))
    return inter / menor          # fraccion del rectangulo MENOR que queda tapada
\end{lstlisting}

El umbral \texttt{0.35} significa que solo se reporta un solape cuando más del 35\,\% del
rectángulo de la palabra menor queda cubierto por otra palabra -- tolera ligeros contactos
de caja (típicos en \texttt{TiKZ} con \texttt{\textbackslash ttk}) y captura colisiones reales
(como la que se corrigió en el diagrama 1 durante el diseño). Con esto el orquestador puede
\emph{verificar sin visión}: reporta ``6 páginas, 0 errores, 1 Overfull, 0 solapes'' en el log
y en el estado.

% ======================================================================
\section{Estado, log y salida JSON}

\subsection{Códigos de salida}

\begin{table}[h!]
\centering
\small
\begin{tabularx}{\linewidth}{@{}lp{2.6cm}>{\raggedright\arraybackslash}X@{}}
\toprule
Código & Origen & Significado \\
\midrule
0 & orquestador & OK: sincronizado o sin cambios. \\
1 & orquestador & Error de cadena (falta archivo, compilación fallida, LLM falló). \\
2 & regenerar & Error de entrada (\texttt{.md}/\texttt{.tex}/\texttt{snapshot} inexistentes). \\
3 & regenerar & Cambio estructural detectado con \texttt{--no-llm}: \texttt{.tex} intacto. \\
4 & regenerar & \texttt{.tex} inexistente: el flujo \emph{sincroniza}, no crea diagramas desde cero. \\
5 & regenerar & Fallo LLM o edits inválidos tras el retry budget (nada escrito). \\
\bottomrule
\end{tabularx}
\end{table}

\subsection{Ejemplos de salidas reales}

Ejecución sin cambios:
\begin{lstlisting}
{"status":"ok","cambio":false,"message":"Sin cambios en el markdown (hash idéntico)."}
\end{lstlisting}

Ejecución con cambio estructural atendido (del test E2E):
\begin{lstlisting}
{"status":"ok","cambio":true,"via":"llm","tex_cambiado":true,
 "secciones":[{"seccion":"2","via":"llm","hunks":1,"edits":1}],
 "edits":1,
 "verificacion":{"status":"ok","paginas":6,"errores":[],
                 "overfull":["Overfull \\hbox (6.79pt too wide) ..."],
                 "solapes":[],"solapes_total":0}}
\end{lstlisting}

El estado persistido tras sincronizar (\texttt{.tmp/faq\_flujo\_sync.json}) contiene los
hashes md/tex, la marca temporal UTC, la vía (determinista/llm/sin\_cambios), las secciones
atendidas, los campos conocidos inyectados y el número total de edits.

% ======================================================================
\section{La directiva (Layer 1)}

La fase 1 del SOP define las entradas requeridas y los 7 pasos; la fase de \emph{edge cases}
es la que hace robusto al flujo. Resumen de los casos límite documentados:

\begin{table}[h!]
\centering
\small
\begin{tabularx}{\linewidth}{@{}p{5cm}X@{}}
\toprule
Caso límite & Protocolo \\
\midrule
Sin snapshot previo (primer arranque) & Se toma línea base: campos deterministas + snapshot; NO se llama al LLM (no hay diff que interpretar). \\
Cambio solo de metadatos (tabla de commits, referencias cruzadas) & Se ignora; el hash queda marcado como atendido sin tocar el \texttt{.tex}. \\
Fallo LLM o edits inválidos & Retry por la cadena de fallback del router (máx 3); si se agota, código 5 sin escribir. \\
El LLM intenta alterar geometría & La validación ``mismo número de \texttt{\textbackslash node[}'' rechaza y aborta. \\
\texttt{--no-llm} con cambio estructural & Código 3; el \texttt{.tex} queda intacto y se reportan las secciones afectadas. \\
\texttt{.tex} inexistente & Código 4 (sincroniza, no crea desde cero). \\
\bottomrule
\end{tabularx}
\end{table}

% ======================================================================
\section{Cómo se validó}

\begin{enumerate}[leftmargin=1.6em]
  \item \textbf{Sanidad estática:} \texttt{py\_compile} de los tres scripts nuevos (se depuró
        un fallo de f-string con barra invertida en el conteo de \texttt{\textbackslash node[}).
  \item \textbf{Pruebas unitarias deterministas} sobre \texttt{\_edits\_validos} /
        \texttt{\_aplicar\_edits}: edit válido preserva el invariante de nodos; \texttt{old}
        no único es rechazado; token estructural prohibido es rechazado; y
        \texttt{aplicar\_campos\_conocidos} inyecta fecha/modelo/N correctos.
  \item \textbf{Línea base real:} un primer arranque crea snapshot + estado; un segundo
        arranque con el \texttt{.md} intacto responde ``hash idéntico'' (sin LLM, sin
        compilación).
  \item \textbf{E2E con LLM real} sobre copias en \texttt{/tmp} (sin tocar el entregable): un
        cambio semántico en la sección 2 fue clasificado como \texttt{llm}, el LLM devolvió
        \texttt{1 edit} quirúrgico correcto (amplió la nota \texttt{N1} con la resolución
        temporal de \texttt{time}, LaTeX válido), la compilación fue exitosa y el verificador
        reportó \texttt{6 páginas, 0 errores, 0 solapes}. El estado real se restauró tras la
        prueba.
\end{enumerate}

\begin{tcolorbox}[cajaConcepto]
\faIcon{check-circle}~El E2E confirmó también el caso \emph{``edits vacíos es correcto''}: el
LLM devolvió \texttt{\{"edits":[]\}} cuando el bloque ya estaba al día, lo que validó el
parseo del JSON vacío y la continuación sin recompilar.
\end{tcolorbox}

% ======================================================================
\section{Uso}

\begin{lstlisting}
python3 flujo_sync_faq_flujo.py                  # sincronización puntual
python3 flujo_sync_faq_flujo.py --watch --interval 60   # bucle en segundo plano
python3 flujo_sync_faq_flujo.py --force          # re-sincronizar ignorando el hash
python3 flujo_sync_faq_flujo.py --no-llm         # solo parte determinista (aborta ante
                                                 #  cambios estructurales, código 3)
python3 flujo_sync_faq_flujo.py --dry-run        # muestra el plan sin escribir/compilar
python3 flujo_sync_faq_flujo.py --critico        # re-traducción LLM en tier opus
python3 execution/regenerar_faq_flujo.py --plan  # solo clasificar (sin LLM, sin escribir)
python3 execution/verificar_pdf.py --pdf build.pdf --log build.log   # guardrail standalone
\end{lstlisting}

\textbf{Coste:} la parte determinista y la detección por hash no consumen créditos. Solo la
re-traducción semántica consume \texttt{OPENROUTER\_API\_KEY} (el orquestador lo avisa antes
de ejecutarla). Con \texttt{--api-backend} no hay opción de backend: \texttt{regenerar} usa
\texttt{openrouter\_chat} por diseño (tier decidido por el enrutador).

% ======================================================================
\section{Limitaciones y trabajo pendiente}

\begin{itemize}[leftmargin=1.4em]
  \item El flujo \textbf{no crea diagramas desde cero} (código 4): la plantilla del diagrama
        es la base; el diseño de un diagrama nuevo sigue siendo tarea del orquestador con
        revisión visual del usuario.
  \item La verificación por bbox detecta solapes de \emph{texto}, no de elementos gráficos
        puros (líneas, flechas); el invariante de nodos y la regla ``solo edita textos''
        protegen la geometría, pero una revisión visual humana del PDF sigue siendo la
        garantía final.
  \item El Overfull cosmético (\texttt{6.79 pt}) pre-existente se reporta y se mantiene;
        no bloquea la cadena.
  \item \emph{Mejora futura sugerida}: usar la telemetría de \texttt{.tmp/routing\_log.jsonl}
        para afinar los umbrales de tokens de la re-traducción (p. ej. cuando una sección
        devuelve consistentemente \texttt{edits:[]}).
\end{itemize}

\begin{tcolorbox}[cajaMejora]
\faIcon{route}~\textbf{En una línea:} editas el FAQ $\rightarrow$ el orquestador detecta el
cambio por hash $\rightarrow$ decide si es determinista o semántico $\rightarrow$ el LLM solo
parchea lo necesario con validación estricta $\rightarrow$ recompila y verifica $\rightarrow$
el PDF queda publicado y el próximo cambio queda armado.
\end{tcolorbox}

\vspace{6pt}
\begin{center}
\small\color{grisTexto}
Documento generado el 12/09/2026. Ficheros fuente: \texttt{flujo\_sync\_faq\_flujo.py},
\texttt{execution/regenerar\_faq\_flujo.py}, \texttt{execution/verificar\_pdf.py},
\texttt{execution/compile\_latex.py} y \texttt{directives/sync\_faq\_a\_flujo.yaml}.
\end{center}

\end{document}